\section{Related Work}
%intro 
 \subsection{Robots in Public and Urban Space}
  Robots have been moving beyond experimental prototypes confined to the laboratory, and gaining greater exposure in public spaces. Jung and Hinds' call for in-the-wild work capable of supporting robust theory marked this turn~\cite{jung2018robots}, and the intervening years have supplied deployments at increasing scale: sidewalk delivery robots now operate commercially~\cite{starship2025}, robotaxis carry a quarter of a million paid rides each week~\cite{waymo2025rides}, and drone delivery is moving from needing waiver to getting routine certification~\cite{faa2025drone}. Studying robots in public therefore matters because deployment changes both the conditions under which robots operate and what counts as successful operation~\cite{jung2018robots}. Outside the laboratory, robots sometimes enter spaces that were not designed around them and encounter people who have not been recruited, trained, or assigned roles in relation to them~\cite{rosenthal2020forgotten}. A robot should navigate and perform its intended task, while people around it must recognize what it is doing~\cite{hetherington2021hey}, decide how to move or act around it~\cite{weinberg2023sharing}, determine whether to help, avoid, or to use it~\cite{weiss2010robots}, and judge whether its presence is appropriate~\cite{bu2025making}.
  
  Field studies have examined how people accommodate, assist, and resist robots in shared spaces. Pedestrians negotiate passage with delivery robots that do not negotiate back~\cite{pelikan2024encountering, yu2024understanding}, compete with them for constrained sidewalk space~\cite{weinberg2023sharing, gehrke2023observed}, and sometimes step in when robots stall~\cite{dobrosovestnova2022little}; for people with mobility disabilities, such encounters could create additional navigation and accessibility challenges~\cite{han2024codesign}. Longer-term deployments similarly show that robots are understood and used differently depending on the people, workflows, and organizations around them~\cite{mutlu2008robots, lee2012ripple, lee2010gracefully}. Other studies in the field also examine breakdowns in public interaction, including cases where people obstruct, damage, or otherwise resist robots: children obstruct and abuse social robots in malls~\cite{brscic2015escaping, nomura2016why}, passers-by push and strike robots in public squares~\cite{salvini2010safe}, and crowds have vandalized autonomous vehicles~\cite{sfstandard2024waymo}. 

  Beyond these direct forms of engagement, robots also affect people who are simply present in the surrounding environment. A related body of work examines such bystanders, who encounter robots without being recruited into interaction, distinguishing them from operators and supervisors~\cite{scholtz2003theory}. Privacy-focused works examine what robots record, who receives those data, and how such information is retained~\cite{rueben2018themes, lutz2019privacy, dietrich2026bystander}. Empirical work further shows that privacy judgments depend on how sensing is framed and on the norms of the setting~\cite{butler2015privacy, denning2014insitu}, consistent with contextual integrity's account of privacy as context-dependent information flow~\cite{nissenbaum2010privacy}.
  
  % Ours is the opposite arrangement: a robot that stays only hours, in six settings, so that what varies is the place rather than the length of the acquaintance.
  
  % A second strand documents what happens when the negotiation fails. Children obstruct and abuse social robots in malls, and asked why, give reasons that are about the robot's perceived agency and about the situation they are standing in~\cite{brscic2015escaping, nomura2016why}; a robot left to move through an Italian city square is pushed, blocked and struck by passers-by~\cite{salvini2010safe}; a Waymo vehicle is set alight by a crowd in San Francisco~\cite{sfstandard2024waymo}. Such incidents are typically narrated either as isolated pathology or as evidence of diffuse public hostility toward automation. Neither reading asks what the robot was taken to \emph{be} at the moment it was attacked, which is the question we think does the explanatory work.
  
  % A third strand concerns the people a public robot records without recruiting. Scholtz's early taxonomy of HRI roles already separates the bystander---present in the robot's space, with no task relationship to it and no way to opt out---from the operator and the supervisor~\cite{scholtz2003theory}, and the privacy literature has since made that position its central problem: what a robot sees, who receives it, and on what terms it may be retained~\cite{kaminski2017averting, rueben2018themes, lutz2019privacy, dietrich2026bystander}. Empirical work here tends to find that concern tracks the framing and the setting rather than the sensor: people asked about a teleoperated home robot's camera respond to how the capability is described~\cite{butler2015privacy}, and bystanders of wearable cameras in public reason about the norms of the specific place they are standing in rather than about recording as such~\cite{denning2014insitu}. Nissenbaum's contextual integrity gives that observation its general form---a flow of information is appropriate or not relative to the norms of the context it occurs in, not in the abstract~\cite{nissenbaum2010privacy}. 
  
Across these lines of work, the challenges of public robot deployment extend beyond technical performance to a broader question of how robots become understandable, acceptable, and workable within existing social environments. Local norms shape what a robot is understood to be and what it is permitted to do~\cite{while2021urban}. The social dimension is therefore not secondary to deployment success, but constitutive of it: even a technically capable robot may fail in practice if people cannot make sense of its purpose or role in the setting. 
Our work builds on this view by examining how place enters into the interpretation of a public robot at the scale of everyday encounters. By deploying the same Trashbot design across sites, we examined how people drew on local social conditions to read the robots' actions and make sense of their role in the setting.
  
  \subsection{Situated Action and Local Context} \label{sec:situated}
A robot's meaning depends on the norms of the place it enters, thus people's reactions to it are \textit{situated}: read off the place as much as off the robot itself. People make sense of a situation by looking at what happened, using cues from the environment and their own experience to build an explanation that seems reasonable~\cite{suchman1987plans, weick2005organizing}. Applied to an encounter with an unfamiliar robot, interpretation draws on what they already know about that place. A robot's onboard camera may be read differently on a university campus than on a commercial street. This is not because the sensor differs, but because the two sites hand the bystander different resources for resolving the same inference~\cite{denning2014insitu}. Site, on this account, is not a backdrop; it supplies the local knowledge the judgment is built from.

Prior work has examined what this local knowledge consists of and how it shapes robot encounters at a given site. Pelikan's argument for \textit{localism} in robot design is the closest statement of the underlying commitment: public space is not a uniform substrate into which robots are inserted, and designing for it requires attending to the specific social organization of specific places~\cite{pelikan2025localism}. Bu et al.'s study of trash barrel robots and Pelikan et al.'s account of street encounters both demonstrate this empirically at a single site, identifying the cultural knowledge, communal common ground, and everyday activity that one place makes available to interpreters~\cite{bu2025making, pelikan2024encountering}. Localism establishes that place shapes interpretation, but it does not yet say what a robot should do to be understood within a given place. 

A related design goal already exists in HRI: \emph{legibility}. Dragan and colleagues define legible motion as motion from which an observer can readily infer a robot's goal, distinguishing it from predictable motion, which matches what an observer expects once that goal is known~\cite{dragan2013legibility}. This idea has since informed work on human-aware and socially aware navigation~\cite{kruse2013human, singamaneni2024survey}. In this literature, legibility is primarily a property of robot behavior: the central question is, \emph{what is this robot about to do?}   

A situated view of interpretation raises a different question: \emph{does the same robot behavior remain equally legible across places?} If people interpret robots using knowledge of the setting around them, then legibility may depend not only on what the robot does, but also on where it does it. Yet little prior work has examined how the interpretation of the same robot changes across local settings. We extend the idea of legibility in two ways. First, we examine the same robot across different sites, allowing us to study how interpretation varies as the surrounding place changes. Second, we turn localism into a design implication. If people across sites ask similar questions---who runs this, is it watching me, does it belong here---then the design challenge is to help the robot answer those shared questions in ways that make sense locally. We call this \emph{local legibility}.

  \subsection{Sensemaking in and beyond HRI}
We use \textit{sensemaking}~\cite{weick2005organizing,pirolli2011introduction} to describe the interpretive work through which people make an unfamiliar situation intelligible and decide how to respond; a process has typically been studied by controlling what the robot does and then asking how people interpret it. This framing shaped three interlocking methodological choices in our study: why we used a Wizard-of-Oz robot, why we situated it through naturalistic observation in a real public space, and why we analyzed interview responses as data about sensemaking.

First, within HRI, Wizard-of-Oz methods have supported a controlled-experiment tradition: by having a hidden operator drive a robot that is perceived as autonomous, researchers can hold its behavior constant while manipulating some other feature of the interaction or the observer, as in studies of the legibility and predictability of robot motion~\cite{dragan2013legibility}, or of how people form mental models and expectations of autonomous systems~\cite{paepcke2010judging, riek2012wizard}. 

Second, we complement this controlled behavioral apparatus with naturalistic observation of bystanders in uncontrolled public space. Previous studies have shown that interpretation occurs even without explicit interaction, as people observe, test, avoid, or otherwise respond to robots encountered in situ~\cite{pelikan2024encountering, sabanovic2006robots}. Because such deployments cannot control what the robot does or when a bystander encounters it, they trade behavioral control for ecological validity. Our approach draws on both traditions: we used the same Trashbot design and Wizard-of-Oz operating arrangement across sites while observing encounters in natural public settings. This allowed us to compare how people read a common service concept across settings while keeping the robot platform and operating arrangement broadly consistent.
  
Third, we treated the post-interaction interviews not as elicited evaluations of the robot but as a further site of sensemaking in their own right. This follows Clark's claim that interpretation depends on what interlocutors take as shared ground~\cite{clark1996using}, and Madsbjerg's argument that people reason through data that is ``thick'' with context rather than abstracted from it~\cite{madsbjerg2017sensemaking}: an interview response is not a transparent readout of what someone thought while encountering the robot, but a further act of sense-making, produced after the fact and shaped by what the person now takes to be relevant. We therefore let the questions we asked stay open, and let the categories we analyzed emerge from what people actually said, rather than fixing in advance which aspects of the robot were worth asking about. The resulting data speak less to how people rated the robot than to what they reached for in order to explain it to themselves and to us--- prior experience with robots, guesses about who was operating it, assumptions about the sites, which is the kind of evidence a closed-ended survey would not have surfaced.
  
  % The term \emph{sensemaking} reaches HCI largely through Weick's account of % organizing, where it names the ongoing work by which people build plausible % accounts of situations that do not interpret % themselves~\cite{weick2005organizing}. HCI has its own lineage for the term--- % sensemaking as the cost-structured work of finding and encoding a representation % that fits a task~\cite{russell1993cost}, and as the gap-bridging people do when a % situation stops cohering~\cite{dervin1998sense}---but Weick's version is the one % that carries the features we need. Three of them bear % directly on what follows. Sensemaking is \emph{retrospective}: people act first % and work out what the action meant afterward, so the account and the conduct are % produced at different moments and need not agree. It is \emph{social}: accounts % are assembled for and with other people, and what settles them is plausibility % within a community rather than correspondence to some fact of the matter. And it % is \emph{ecological}: the materials people reason with are furnished by the % situation they are standing in, not retrieved whole from memory. Weick's % sensemaker is not a mind solving a puzzle but a person acting into an equivocal % world and talking about it as they go.
  % HRI has adopted the vocabulary while narrowing the construct. In much of the % field's usage, sensemaking names something a single person does internally during % an encounter: an inferential process running from cues on the robot's surface---its % form factor, its motion, its markings---to a conclusion about what the machine is % and what it wants. The unit of analysis is the individual; the timescale is the % encounter; the outcome is a mental model. This framing has been productive, and % the closest prior study of trash barrel robots in public space works within % it~\cite{bu2025making}, as does the broader literature on how bystanders read % robots they meet on the street~\cite{pelikan2024encountering}. But it quietly % drops two of Weick's three features. The social production of the account % disappears into individual cognition, and the environment shrinks to the robot % itself---a source of stimuli rather than a source of interpretive resources.
  % Our contribution is in part a recovery. We argue that the environment people % recruit when they make sense of a public robot is not the robot but the % \emph{place}: its history, its reputation, the work that maintains it, the bodies % that occupy it, the institutions understood to be operating in it. Restoring the % setting to the analysis changes what counts as an explanation of cross-community % difference. If sensemaking is an individual cognitive act, differences between % neighborhoods must be explained by differences between the people in them---by % attitudes, familiarity, or demographics that individuals carry with them. If it % is a situated practice, the same explanation is available without that % commitment: the practice is shared, and the difference lies in what each place % puts within reach. % \note{Verify the characterization of \cite{bu2025making} and % \cite{pelikan2024encountering} against the actual papers before submission --- % this paragraph asserts what framing they adopt, and a reviewer who knows either % work will check. If Pelikan et al.\ are already closer to a situated reading than % implied here, soften to cite them as a partial precedent rather than an instance % of the narrow framing.}
 